% Irodov, user-provided PDF, original pp.41--43, PDF pp.50--52.
% Problem 264's incline slopes DOWN as drawn in Fig. 1.67, not uphill.

\exproblem{261}{41}
质量为 $m_1$ 的木板置于光滑水平面上，
质量为 $m_2$ 的均匀实心球置于板上。
对木板施加恒定水平力 $F$，
假设球与板之间始终无相对滑动。
求木板与球心的水平加速度 $a_1,a_2$。
\exsolution
设板向右加速，静摩擦力 $f$ 使球心向右加速并
使球绕自身中心逆时针自转。对球有
\[
m_2a_2=f,\qquad I_C\beta=fR,\qquad
I_C=\frac25m_2R^2.
\]
球与板接触处无相对切向加速度，
因此
$a_1=a_2+R\beta$。
消去 $\beta,f$：
\[
a_1=a_2+\frac{fR^2}{I_C}
=a_2+\frac52a_2=\frac72a_2.
\]
木板的水平运动方程
$F-f=m_1a_1$，代入 $f=m_2a_2=(2/7)m_2a_1$：
\[
\boxed{a_1=\frac{F}{m_1+\frac27m_2}},\qquad
\boxed{a_2=\frac27a_1}.
\]
\excheck 球心比木板加速得慢，
但无滑动条件约束的是球体接触点的速度，
不是两者质心速度相等。

\exproblem{262}{41--42}
均匀实心圆柱质量为 $m$、半径为 $R$，
最初仅绕自身轴以角速度 $\omega_0$ 自转、质心静止，
随后被放到水平粗糙面上释放。
摩擦系数为 $k$。求：
(a) 从释放到开始纯滚动的滑动持续时间；
(b) 滑动摩擦对圆柱所做\emph{总功}。
\exsolution
滑动期间地面摩擦力大小恒为 $f=kmg$。
其使圆柱获得质心速度，$a_C=kg$，
同时施加阻转力矩，
\[
|\beta|=\frac{fR}{I_C}=\frac{2kg}{R},
\qquad I_C=\frac12mR^2.
\]
以摩擦使质心前进的方向为正，
\[
v_C(t)=kgt,\qquad
\omega(t)=\omega_0-\frac{2kg}{R}t.
\]
首次纯滚动条件为 $v_C=R\omega$：
\[
kgt=\omega_0R-2kgt
\Longrightarrow
\boxed{t_{\mathrm{slide}}=\frac{\omega_0R}{3kg}}.
\]
此时 $v_f=\omega_0R/3,\ \omega_f=\omega_0/3$。
摩擦对整个圆柱的功等于总动能变化：
\[
W_f=\left[\frac12mv_f^2+\frac12I_C\omega_f^2\right]
-\frac12I_C\omega_0^2
=\left(\frac1{18}+\frac1{36}-\frac14\right)
mR^2\omega_0^2.
\]
故
\[
\boxed{W_f=-\frac16mR^2\omega_0^2}.
\]
\excheck 在滑动阶段存在机械能耗散；
接触点速率由不为零逐渐降至零，之后进入纯滚动。

\exproblem{263}{42}
一颗均匀实心小球半径为 $r$，
在半径为 $R$ 的固定大球表面从顶端几乎静止开始
无滑动滚下。
求小球恰离开大球表面瞬间的自转角速度大小。
\exsolution
小球质心绕大球心移动的圆半径为 $L=R+r$。
设质心的径矢与竖直向上方向夹角为 $\theta$，
质心速率为 $v$，小球自转角速度
$\omega=v/r$。
由能量守恒：
\[
mgL(1-\cos\theta)
=\frac12mv^2+\frac12\left(\frac25mr^2\right)\frac{v^2}{r^2}
=\frac7{10}mv^2.
\]
故
$v^2=\frac{10}{7}gL(1-\cos\theta)$。
径向向心方向列牛顿第二定律：
\[
mg\cos\theta-N=\frac{mv^2}{L}.
\]
离面瞬间 $N=0$，联立得到
$\cos\theta=10/17$，
$v^2=(10/17)gL$。
从而
\[
\boxed{\omega_{\mathrm{detach}}=
\frac vr=\sqrt{\frac{10g(R+r)}{17r^2}}}.
\]
\excheck 半径为 $R+r$ 的圆是\emph{球心轨迹}，
滚动约束的自转半径则为 $r$，不能把二者混用。

\exproblem{264}{42}
半径 $R=15\,\mathrm{cm}$ 的均匀实心圆柱
原来以速率 $v_0$ 在水平面上向右纯滚动，
其前方路面在拐角处\emph{向下}折成
倾角 $\alpha=30^\circ$ 的斜面（原图1.67）。
求使圆柱进入斜面段时\emph{不腾空跳离拐角}的
最大 $v_0$，并假设运动中无滑动。
\exsolution
图1.67的斜面向右\emph{下倾}。
圆柱过拐角的短阶段以拐角 $O$ 为瞬时接触点，
圆心沿半径 $R$ 的圆弧从竖直上方转过
$\theta\in[0,\alpha]$。
圆柱绕接触点的转动惯量是
$I_O=I_C+mR^2=\frac32mR^2$。
由于纯滚动 $v=R\dot\theta$，
其总动能 $K=\frac12I_O\dot\theta^2
=\frac34mv^2$。
圆心在下坡阶段降低
$R(1-\cos\theta)$，能量守恒：
\[
\frac34mv^2
=\frac34mv_0^2+mgR(1-\cos\theta),
\quad
v^2=v_0^2+\frac43gR(1-\cos\theta).
\]
沿指向拐角的径向列式：
$mg\cos\theta-N=mv^2/R$，
其中 $N$ 为拐角向外的支持力。
要一直保持接触须 $N\ge0$，即
\[
\frac{N}{mg}=\frac73\cos\theta-\frac43
-\frac{v_0^2}{gR}\ge0.
\]
由于 $\cos\theta$ 在区间中递减，
最严条件发生在 $\theta=\alpha$，于是
\[
\boxed{v_{0,\max}=
\sqrt{\frac{gR}{3}(7\cos\alpha-4)}
\simeq1.00\,\mathrm{m/s}}.
\]
\excheck 根号内须为非负，即
$\cos\alpha\ge4/7$；本题 $\alpha=30^\circ$ 满足。
如果误把原图的向下斜坡当成向上斜坡，
得到的离面条件将完全不同。

\exproblem{265}{42}
一根均匀薄刚性圆环质量为 $m$、半径为 $R$，
其内缘固定着质量同样为 $m$ 的小物体 $A$。
圆环在水平面上无滑动滚动。
每当 $A$ 到达环内最低位置时，
圆环中心速率为 $v_0$（图1.68）。
求圆环始终不跳离地面的 $v_0$ 条件。
\exsolution
令 $\theta=0$ 表示小物体位于环心正下方，
圆环滚动时其中心速率为 $v=R\dot\theta$。
环本身的平动与转动动能之和为 $mv^2$；
小物体绕地面接触点作摆线运动，其速率平方为
$v_A^2=2v^2(1-\cos\theta)$，
故全系统动能
\[
K=mv^2+\frac12m v_A^2
=mv^2(2-\cos\theta).
\]
小物体重力势能相对最低点的增量为
$mgR(1-\cos\theta)$。能量守恒：
\[
v^2(2-\cos\theta)+gR(1-\cos\theta)=v_0^2.
\]
环和小物体总质心高度为
$y_C=R-\frac12R\cos\theta$，
质量合计 $2m$。
地面对系统的支持力
$N=2mg+2m\,\ddot y_C$；
代入上式与其关于时间的导数可确定
$N(\theta)$。在足以完成整圈的运动中，
地面支持力的最小值出现在
$A$ 位于最高位置 $\theta=\pi$。
此时 $\dot\theta^2=v^2/R^2$、
$\ddot y_C=-R\dot\theta^2/2$，于是
$N_{\mathrm{top}}=2mg-mv_{\mathrm{top}}^2/R$。
能量式在 $\theta=\pi$ 给出
$v_{\mathrm{top}}^2=(v_0^2-2gR)/3$。
要求 $N_{\mathrm{top}}\ge0$，
即
\[
\frac{v_0^2-2gR}{3}\le2gR
\quad\Longrightarrow\quad
\boxed{v_0\le\sqrt{8gR}}.
\]
\excheck 若 $v_0$ 太小，小物体不会到达圆环顶端，
但系统仍可能往复滚动而不跳起；
式中给的是无跳离地面的上限，原书亦如此给出。

\exproblem{266}{42}
一辆拖拉机以相对于地面恒速 $v$ 前进，
整条履带总质量为 $m$，如图1.69。
在无滑动滚动条件下，求整条履带相对地面的总动能。
\exsolution
在理想履带模型中，上方直线段沿前进方向
相对地面速率为 $2v$；底部接触地面的直线段
瞬时速率为零。忽略弯曲端部的质量，
两段履带质量分别为 $m/2$。
因此
\[
K=\frac12\left(\frac m2\right)(2v)^2
+\frac12\left(\frac m2\right)0^2
=\boxed{mv^2}.
\]
\excheck 若端部弯曲段质量不可忽略，应将沿履带各质元
的实际速度积分，此处采用原题图的理想化直线段近似。

\exproblem{267}{42}
质量为 $m$、半径为 $r$ 的均匀实心球
绕水平杆 $OA$ 转动并在水平面上无滑动滚动
（原图1.70）。球心沿距竖直旋转中心 $O$
为 $R$ 的圆轨道以速率 $v$ 前进。
求球的总动能。
\exsolution
球心平动动能为 $mv^2/2$。
运动中球同时绕水平径向轴滚动，角速度大小
$\omega_r=v/r$；球心的轨道绕竖直方向以
$\Omega=v/R$ 转动。两转动角速度矢量互相垂直。
实心球各直径轴惯量相同，均为
$I_C=\frac25mr^2$，故球自转动能为
\[
K_{\rm rot}=\frac12I_C(\omega_r^2+\Omega^2)
=\frac15m v^2\left(1+\frac{r^2}{R^2}\right).
\]
总动能
\[
\boxed{K=\frac7{10}mv^2\left(1+\frac27\frac{r^2}{R^2}\right)}.
\]
\excheck 若遗漏球心绕中心 $O$ 的方位转动角速度
$\Omega$，将丢失与 $r^2/R^2$ 成比例的附加动能。

\exproblem{268}{42--43}
证明：一个质量为 $m$、质心相对于固定旋转轴
的垂直位置矢量为 $\bm R_C$ 的刚体，
在以恒角速度 $\bm\omega$ 旋转的参考系中，
全部质元所受的离心惯性力合力和科里奥利惯性力合力分别为
\[
\bm F_{\rm cf}=m\omega^2\bm R_C,\qquad
\bm F_{\rm Cor}=2m(\bm v_C'\times\bm\omega),
\]
其中 $\bm v_C'$ 为质心在旋转系中的速度。
\exsolution
把刚体分成质元 $\mathrm dm$。
各质元离心惯性力
$\mathrm d\bm F_{\rm cf}
=-\bm\omega\times(\bm\omega\times\bm r_\perp)
\,\mathrm dm$。
由于 $\bm r_\perp$ 垂直于 $\bm\omega$，
双叉乘恒等式给出
$-\bm\omega\times(\bm\omega\times\bm r_\perp)
=\omega^2\bm r_\perp$。
对全体质量元积分：
\[
\bm F_{\rm cf}
=\omega^2\int\bm r_\perp\,\mathrm dm
=m\omega^2\bm R_C.
\]
各质元科里奥利惯性力为
$\mathrm d\bm F_{\rm Cor}
=-2\bm\omega\times\bm v_i'\,\mathrm dm$。
利用积分与叉积的线性性及质心速度定义：
\[
\bm F_{\rm Cor}=-2\bm\omega\times
\int\bm v_i'\,\mathrm dm
=-2m\bm\omega\times\bm v_C'
=\boxed{2m(\bm v_C'\times\bm\omega)}.
\]
\excheck 这些是\emph{合力}恒等式，不表示各质元受惯性力
均相同；求惯性力矩仍需考虑物体内部质量分布。

\exproblem{269}{43}
长为 $\ell$、质量为 $m$ 的均匀细杆中点 $C$
固定于旋转轴 $OO'$ 上，杆与轴夹角为 $\theta$
（原图1.71）。轴以恒定角速度 $\omega$ 转动。
求在随轴旋转的参考系中，全部离心惯性力
对杆中点 $C$ 的\emph{合力矩大小}。
\exsolution
令杆上质元相对 $C$ 的位置为
$\bm r=s\bm e$，其中
$s\in[-\ell/2,\ell/2]$、$\bm e$ 是杆轴单位矢量，
与竖直旋转轴夹角为 $\theta$。
离心惯性力为
$\mathrm d\bm F=\omega^2(s\sin\theta)
\bm e_\perp\,\mathrm dm$，
其中 $\bm e_\perp$ 是杆方向的水平投影的单位矢量。
每个微元关于 $C$ 的力矩模为
\[
\mathrm dN
=\omega^2s^2\sin\theta\cos\theta\,\mathrm dm,
\]
所有微元力矩同向。因均匀细杆
$\int s^2\,\mathrm dm=m\ell^2/12$，
积分得
\[
\boxed{N=\frac{m\omega^2\ell^2}{12}
\sin\theta\cos\theta
=\frac{m\omega^2\ell^2}{24}\sin2\theta}.
\]
\excheck 这也可由刚体惯量张量的陀螺力矩公式得到。
本版书末第269题印作 $\frac12m\omega^2\ell^2\sin2\theta$，
与上述积分及同一原书第282(c)题的
$\frac1{24}m\omega^2\ell^2\sin2\theta$ 不一致；
本稿保留可复核的连续体积分结论并标为勘误。

\exproblem{270}{43}
一圆锥摆由长为 $\ell$、质量为 $m$ 的均匀细杆构成，
上端通过铰链与固定点连接。
细杆绕竖直轴以恒角速度 $\omega$ 匀速旋转，
求其与竖直方向的平衡夹角 $\theta$。
\exsolution
在随杆匀速旋转的参考系中，
杆处于相对静止状态。
选铰点为力矩中心，铰接反力无力矩。
重力作用在杆中点，其使杆向竖直方向回落的力矩大小
\[
N_g=mg\frac{\ell}{2}\sin\theta.
\]
取距铰点沿杆距离为 $s$ 的质量微元，
其到竖直轴距离为 $s\sin\theta$，
离心力大小 $\omega^2s\sin\theta\,\mathrm dm$，
对铰点的力矩大小是
$\omega^2s^2\sin\theta\cos\theta\,\mathrm dm$。
积分且 $\int s^2\,\mathrm dm=m\ell^2/3$：
\[
N_{\rm cf}=\frac13m\omega^2\ell^2
\sin\theta\cos\theta.
\]
令力矩平衡 $N_g=N_{\rm cf}$，
对非竖直解约去 $\sin\theta$：
\[
\boxed{\cos\theta=\frac{3g}{2\omega^2\ell}}.
\]
\excheck 要存在非竖直圆锥摆，
必须 $\omega^2\ell\ge3g/2$；
低于该阈值时仅有 $\theta=0$ 的竖直平衡。
